\section{Method}
\label{sec:method}




We present \textbf{NESS}, a framework for learning under missing node features.
NESS is deliberately modular: a feature-propagation prefill and a message-passing
encoder produce node representations, onto which a set of \emph{pluggable}
self-supervised objectives is layered. This modularity is intentional. Our study
of which objectives help under missingness (Section~\ref{sec:experiments})
is enabled by treating the auxiliary objective as an interchangeable component
rather than a fixed part of the architecture.




\begin{figure}
    \centering
    \includegraphics[width=0.85\linewidth]{figures/TGINA-method.pdf}
    \caption{The pipeline of NESS. We first apply Feature Propagation to prefill missing nodes. We then optimize the contrastive and self-supervised objectives. Green boxes indicate loss objectives.}
    \label{fig:method}
\vspace{-15pt}
\end{figure}



\subsection{Problem Formulation and Notation}
\label{sec:method-problem}

Let $G=(V,E,\mathbf{X})$ be an undirected attributed graph with $n=|V|$ nodes,
edge set $E$, adjacency matrix $\mathbf{A}\in\{0,1\}^{n\times n}$, and feature
matrix $\mathbf{X}\in\mathbb{R}^{n\times d}$. Each node $v_i$ has a label
$y_i\in\{1,\dots,C\}$. We denote by $\tilde{\mathbf{A}}=\mathbf{D}^{-1/2}(\mathbf{A}
+\mathbf{I})\mathbf{D}^{-1/2}$ the symmetrically normalized adjacency with
self-loops, where $\mathbf{D}$ is the degree matrix.

Under \emph{attribute-missing} conditions, the node set partitions into an
\textbf{observable} set $V_o$ and a \textbf{missing} set $V_m$, with $V_o\cup V_m
=V$ and $V_o\cap V_m=\emptyset$. For a missing node $v_i\in V_m$, the entire
feature vector is unobserved; we represent the observed input by zero-filling,
\begin{equation}
    \label{eq:mask}
    \bar{\mathbf{x}}_i =
    \begin{cases}
        \mathbf{x}_i, & v_i \in V_o, \\
        \mathbf{0},   & v_i \in V_m,
    \end{cases}
\end{equation}
giving the masked feature matrix $\bar{\mathbf{X}}$. Following a completely random missing pattern, $V_m$ is sampled uniformly at a rate $\rho=|V_m|/n$.
The goal is \emph{not} to reconstruct $\mathbf{X}$ but to learn node
representations $\mathbf{Z}\in\mathbb{R}^{n\times h}$ that support accurate
classification of the missing nodes $V_m$, measured by macro-F1.

\subsection{Overview}
\label{sec:method-overview}

NESS comprises four stages (Figure~\ref{fig:method}): (1)~a \emph{feature-propagation prefill} that fills missing rows of $\bar{\mathbf{X}}$ with long-range
structural signal; (2)~a message-passing \emph{encoder} that maps the prefilled
features to representations $\mathbf{Z}$; (3)~a \emph{classification head} trained
on observable labels; and (4)~a set of \emph{self-supervised objectives} that
shape $\mathbf{Z}$. The total training loss combines the supervised and
self-supervised terms. We describe each stage below.

\subsection{Feature-Propagation Prefill}
\label{sec:method-prefill}

A shallow message-passing encoder starves when a node's neighborhood is largely
empty, which is precisely the regime induced by high missingness. To supply
missing nodes with signal before encoding, we prefill $\bar{\mathbf{X}}$ using
Feature Propagation~\cite{rossi2022unreasonable}: features are diffused over the graph while observed
rows are reset to their true values at every step. Starting from
$\mathbf{H}^{(0)}=\bar{\mathbf{X}}$,
\begin{equation}
    \label{eq:fp}
    \mathbf{H}^{(t+1)} = \tilde{\mathbf{A}}\,\mathbf{H}^{(t)},
    \qquad
    \mathbf{H}^{(t+1)}_{i,:} = \mathbf{x}_i \;\; \forall\, v_i \in V_o,
\end{equation}
iterated for $T$ steps (we use $T=40$). The prefilled matrix
$\hat{\mathbf{X}}=\mathbf{H}^{(T)}$ replaces $\bar{\mathbf{X}}$ as the encoder
input. The reset in \eqref{eq:fp} enforces the observed values as boundary
conditions, so the diffusion solves a Dirichlet energy minimization that
propagates observed features into missing regions. Prefill is a one-time
preprocessing step of cost $O(|E|\,d\,T)$.

\subsection{Encoder and Views}
\label{sec:method-encoder}

The encoder $f_\theta$ is an $L$-layer GraphSAGE~\cite{hamilton2017inductive} network. At layer
$\ell$, node $v_i$ updates its representation by combining a transform of its own
state with an aggregate of its neighbors $\mathcal{N}(i)$,
\begin{equation}
    \label{eq:sage}
    \mathbf{h}_i^{(\ell)} = \sigma\!\Big(
        \mathbf{W}_{\text{self}}^{(\ell)}\mathbf{h}_i^{(\ell-1)}
        + \mathbf{W}_{\text{nbr}}^{(\ell)}\,
        \mathrm{mean}_{j\in\mathcal{N}(i)}\,\mathbf{h}_j^{(\ell-1)}
    \Big),
\end{equation}
with $\mathbf{h}_i^{(0)}=\hat{\mathbf{x}}_i$ and $\sigma$ a nonlinearity. The
separate self-transform $\mathbf{W}_{\text{self}}$ preserves each node's own
(prefilled) signal through depth, which we find important under missingness
(Section~\ref{sec:experiments}).

To support view-based objectives, we form two stochastic views by independent
random edge-masking of $G$, encode each, and fuse them into the final
representation,
\begin{equation}
    \label{eq:views}
    \mathbf{Z}^{(1)} = f_\theta(\hat{\mathbf{X}}, G^{(1)}),\quad
    \mathbf{Z}^{(2)} = f_\theta(\hat{\mathbf{X}}, G^{(2)}),\quad
    \mathbf{Z} = \tfrac{1}{2}\big(\mathbf{Z}^{(1)}+\mathbf{Z}^{(2)}\big),
\end{equation}
where $G^{(1)},G^{(2)}$ are edge-masked variants of $G$, and we denote by $E_m$ the edges removed to form $G^{(1)}$, which serve as targets for the edge-reconstruction objective below. The fused $\mathbf{Z}$ is
used by the classification head and all objectives. The view-generation method is itself a modular choice; we adopt independent edge-masking, which we found effective in our setting.

\subsection{Learning Objectives}
\label{sec:method-objectives}

\textbf{Classification:}
A head $g_\phi$ maps representations to class logits, trained with cross-entropy on
\emph{observable} nodes only (no label leakage to $V_m$),
\begin{equation}
    \label{eq:cls}
    \mathcal{L}_{\text{cls}}
    = \frac{1}{|V_o|}\sum_{v_i\in V_o}
      \mathrm{CE}\big(g_\phi(\mathbf{z}_i),\, y_i\big).
\end{equation}

\textbf{Edge reconstruction:}
A structural objective encourages $\mathbf{Z}$ to preserve connectivity. The positives are the edges held out when forming view $G^{(1)}$ in \eqref{eq:views}; using an
edge decoder $\psi$, these are contrasted against sampled
negatives $E^-$, a set of randomly drawn non-adjacent node pairs,
\begin{equation}
    \label{eq:edge}
    \mathcal{L}_{\text{edge}}
    = -\!\!\sum_{(i,j)\in E_m}\!\!\log \sigma\big(\psi(\mathbf{z}_i,\mathbf{z}_j)\big)
      -\!\!\sum_{(i,j)\in E^-}\!\!\log\big(1-\sigma(\psi(\mathbf{z}_i,\mathbf{z}_j))\big),
\end{equation}
where $E_m$ is the set of masked edges.

\textbf{Contrastive (view agreement):}
We align the two views with a redundancy-reduction (Barlow-Twins~\cite{zbontar2021barlow})
objective on their cross-correlation matrix $\mathcal{C}$, computed over the batch
dimension of the projected views,
\begin{equation}
    \label{eq:con}
    \mathcal{L}_{\text{con}}
    = \sum_{k}(1-\mathcal{C}_{kk})^2
      + \lambda \sum_{k}\sum_{k'\neq k}\mathcal{C}_{kk'}^2 .
\end{equation}


\textbf{Self-supervised auxiliary objective:} The distinctive component of NESS is a pluggable self-supervised objective $\mathcal{L}_{\text{ssl}}$ that shapes $\mathbf{Z}$ using structure- or label-derived targets requiring no missing-node features. Rather than committing to a single objective, we treat $\mathcal{L}_{\text{ssl}}$ as an interchangeable component and study a family of candidates spanning distinct design axes: \emph{feature-neighborhood} (predict the neighborhood mean or spread), \emph{feature-reconstruction} (masked-feature recovery), \emph{structural} (multi-hop reachability, distance ranking, anchor-distance to landmarks), and \emph{label-based} (neighbor label-histogram, semi-supervised). Each is defined formally in Appendix~\ref{app:objectives}; their comparative analysis, which objectives help and when, is the subject of Section~\ref{sec:experiments}.

Our reported model instantiates $\mathcal{L}_{\text{ssl}}$ with two objectives. The label-histogram objective predicts, for each observable node, the class distribution $\mathbf{p}_i$ of its observable neighbors, using a head $h_{\text{hist}}$ trained with a KL divergence,
\begin{equation}
    \label{eq:hist}
    \mathcal{L}_{\text{hist}} = \frac{1}{|V_o|}\sum_{v_i\in V_o}
      \mathrm{KL}\big(\mathbf{p}_i \,\|\, \mathrm{softmax}(h_{\text{hist}}(\mathbf{z}_i))\big).
\end{equation}

The label-histogram target is computed exclusively from the labels of observable (training) nodes; validation and test labels never enter any target, so the objective introduces no leakage of evaluation labels. The target for node $i$ is an aggregate of its graph neighbors' labels: the adjacency carries no self-loops, so a node's own label does not appear in its one-hop histogram. Because the objective reuses only labels already available to the supervised loss, it adds no information beyond the standard training signal, and is best understood as a semi-supervised auxiliary target rather than a source of leakage.


The multi-hop reachability (path) objective samples random walks to obtain node pairs $(s,e)$ reachable within a few hops as positives and distant pairs $(s,e^-)$ as negatives, and scores them with a head $h_{\text{path}}$,
\begin{equation}
    \label{eq:path}
    \mathcal{L}_{\text{path}} = \tfrac{1}{2}\Big[
        \mathrm{BCE}\big(h_{\text{path}}(\mathbf{z}_s\!\odot\!\mathbf{z}_e),\,1\big)
        + \mathrm{BCE}\big(h_{\text{path}}(\mathbf{z}_s\!\odot\!\mathbf{z}_{e^-}),\,0\big)
    \Big].
\end{equation}
Here $\mathcal{L}_{\text{ssl}}=\mathcal{L}_{\text{hist}}+\mathcal{L}_{\text{path}}$.


\textbf{Total objective:}
The training loss is a weighted combination,
\begin{equation}
    \label{eq:total}
    \mathcal{L}
    = w_{\text{cls}}\,\mathcal{L}_{\text{cls}}
    + w_{\text{edge}}\,\mathcal{L}_{\text{edge}}
    + w_{\text{con}}\,\mathcal{L}_{\text{con}}
    + w_{\text{ssl}}\,\mathcal{L}_{\text{ssl}},
\end{equation}
with nonnegative weights $w_{\cdot}$. Setting $w_{\text{ssl}}=0$ recovers a strong
no-SSL baseline used throughout our analysis; it still retains the classification, edge-reconstruction, and contrastive terms, and drops only the auxiliary self-supervised objective.

\subsection{Training}
\label{sec:method-training}
For large graphs, we partition $G$ into clusters following~\cite{chiang2019cluster} and train
full-batch within each cluster; small graphs are trained full-batch directly.
Self-supervised targets that depend on non-local information (e.g.\ landmark distances, or
label histograms aggregated over a multi-hop neighborhood) are computed once on the full graph and indexed per cluster.
The per-epoch cost is dominated by the
encoder forward/backward pass, $O(|E|\,h + n\,h^2)$ per cluster, matching standard
message-passing GNNs. Full hyperparameter and optimization settings are given in Section~\ref{sec:setup}.
